\section{Wnioski i pytania otwarte}

\begin{stwierdzenie}[ZWERYFIKOWANE]\label{st:gora}
Istnieje automorfizm $F\in\Aut(\Q^2)$ stopnia $6$ z $F(p_i)=q_i$, $i=1,\dots,5$, np.
$F=A\circ u_g\circ t_f\circ B$ z $f=\tfrac83x^2$, $g=-\tfrac5{96}y^3-\tfrac1{24}y^2$
(plik \texttt{praca/iter02/zad3/rozwiazanie\_st6\_ladne.json}). Weryfikator przyjął
łącznie 4 różne pliki stopnia 6 (multistopnie $(2,3)$ i $(3,2)$).
\end{stwierdzenie}

\begin{lemat}[OBLICZENIOWE, trzy niezależne skrypty]\label{lem:funkc}
Jeśli $l,l'$ są funkcjonałami afinicznymi nad $\Q$ i $l(p_i)=l'(q_i)$ dla $i=1,\dots,5$, to
$l=l'$ jest stałą (macierz $[p_x,p_y,1,-q_x,-q_y,-1]$ ma rząd 5).
\end{lemat}

\begin{lemat}[OBLICZENIOWE, dwie metody]\label{lem:22}
Nie istnieje $F=A\circ u_g\circ t_f\circ L$ z $A,L\in\mathrm{Aff}(\Q)$, $\deg f=\deg g=2$,
realizujący bijekcję. Wielomian eliminacyjny dla $s=\alpha_1^3g_2$ rozkłada się na czynnik
kwadratowy (ciało $\Q(\sqrt{94})$) i kubiczny, bez pierwiastków wymiernych. Nad odpowiednimi
ciałami liczbowymi takie automorfizmy stopnia 4 istnieją.
\end{lemat}

\begin{twierdzenie}[UDOWODNIONE modulo tw.~\ref{tw:jvdk} i lematy \ref{lem:funkc}, \ref{lem:22}]\label{tw:dol}
Każdy automorfizm wielomianowy $F$ płaszczyzny $\Q^2$ z $F(p_i)=q_i$ ma $\deg F\ge6$.
W konsekwencji, wraz ze stwierdzeniem~\ref{st:gora}, \textbf{najmniejszy stopień takiego
automorfizmu wynosi 6}.
\end{twierdzenie}
\begin{proof}
Słowa długości $k=0$: $F$ afiniczne, więc $l=x$, $l'=$ pierwsza składowa $F^{-1}$ są niestałe
i zgodne -- sprzeczność z lematem~\ref{lem:funkc}. Dla $k=1$, $F=A\circ T\circ L$ z $T\in J$
zachowującym pierwszą współrzędną z dokładnością do $\alpha\,\mathrm{pr}_1+c$, funkcjonały
$l=\alpha\,\mathrm{pr}_1\circ L+c$, $l'=\mathrm{pr}_1\circ A^{-1}$ są niestałe i zgodne --
sprzeczność. Zatem $k\ge2$ i z lematu~\ref{lem:postac} $\deg F$ jest iloczynem co najmniej dwóch
liczb $\ge2$. Stąd $\deg F\notin\{1,2,3,5\}$, a $\deg F=4$ wymaga $k=2$, $\deg f_1=\deg f_2=2$,
czyli $F=A\circ u_g\circ t_f\circ L$ z $\deg f=\deg g=2$, co wyklucza lemat~\ref{lem:22}.
\end{proof}

\begin{uwaga}
Status twierdzenia~\ref{tw:dol} jest mieszany: część rozumowania jest udowodniona, ale zależy
od cytowanego (nieweryfikowanego tu) twierdzenia Junga--van der Kulka oraz od dwóch rachunków
komputerowych w Sage/Singularze i sympy. Wyniki te potwierdzono niezależnymi skryptami, nie
zostały jednak sformalizowane. Lemat~\ref{lem:22} pokazuje też, że ograniczenie dolne jest
,,arytmetyczne'': nad ciałem $\Q(\sqrt{94})$ stopień 4 byłby osiągalny.
\end{uwaga}

\subsection*{Pytania otwarte}
\begin{enumerate}
\item Czy istnieje bezkomputerowy (ręczny) dowód lematu~\ref{lem:22}, np. przez jawny
opis rozwiązań nad $\overline{\Q}$?
\item Jak opisać zbiór wszystkich automorfizmów stopnia 6 realizujących bijekcję? Przegląd
kierunków $|\cdot|\le3$ dał tylko 6 par $(a,e)$ (wynik obliczeniowy); czy przy kierunkach
wymiernych rodzina jest skończona z dokładnością do trywialnej swobody?
\item Czy dla ogólnych (,,generycznych'') danych $m=5$ najmniejszy stopień nad $\Q$ wynosi
zawsze 6, a nad $\overline{\Q}$ -- 4? (HIPOTEZA, niebadana.)
\item Jak zachowuje się najmniejszy stopień dla większych $m$ w porównaniu z ograniczeniem
$(m-1)^2$?
\end{enumerate}
